
\chapter{Die kleinen Helfer -- Vokale}

Arabisch schreibt normalerweise nur die Konsonanten und langen Vokale.
Die kurzen Vokale (a, i, u) werden meist \emph{nicht} geschrieben -- du musst
sie lernen wie Melodien. In Lehrbüchern (und hier) setzen wir sie aber als
kleine Hilfszeichen über oder unter den Buchstaben.

\begin{center}
\renewcommand{\arraystretch}{2}
\begin{tabular}{cccl}
\textbf{Name} & \textbf{Lage} & \textbf{Beispiel} & \textbf{Klingt wie} \\
\hline
Fatha & Strich oben & \textarabic{\Large بَ} (ba) & kurzes a \\
Kasra & Strich unten & \textarabic{\Large بِ} (bi) & kurzes i \\
Damma & Komma oben & \textarabic{\Large بُ} (bu) & kurzes u \\
Sukun & Kringel oben & \textarabic{\Large بْ} (b) & gar kein Vokal \\
Shadda & kleines w oben & \textarabic{\Large بّ} (bb) & doppelter Konsonant \\
\end{tabular}
\end{center}

\begin{lachbox}
Fatha sieht aus wie ein kleiner, gefallener Schnurrbart über dem
Buchstaben. Kasra ist derselbe Schnurrbart, nur er ist heruntergefallen
und liegt jetzt unten. Damma ist ein winziges Ohr, das nach oben lauscht.
\end{lachbox}

\section*{Die drei langen Vokale}

Lange Vokale werden dagegen tatsächlich als eigene Buchstaben geschrieben:
\arbig{ا} für langes \emph{aa}, \arbig{و} für langes \emph{uu}, \arbig{ي}
für langes \emph{ii}.

\begin{center}
\begin{tabular}{ll}
langes \textit{aa}: & \textarabic{\Large كتاب} \quad kit\'ab (Buch) \\
langes \textit{uu}: & \textarabic{\Large دور} \quad duur (Rollen/Runden) \\
langes \textit{ii}: & \textarabic{\Large كبير} \quad kab\'ir (groß) \\
\end{tabular}
\end{center}

\begin{tippbox}
Ein Wort, das du garantiert brauchen wirst: \textarabic{\Large حبيبي} --
\textit{Hab\'ibi} -- ,,mein Schatz/Liebling``. Langes \textit{ii} mitten
im Wort. Wenn Mohamed dich so nennt, weißt du jetzt auch, warum es sich
so lang anhört.
\end{tippbox}

\section*{Kurzer Vergleich}

\begin{center}
\begin{tabular}{lll}
\textarabic{\Large كتاب} kit\'ab & (ein) Buch & mit langem a \\
\textarabic{\Large كُتُب} kutub & Bücher (Plural!) & mit kurzem u, u \\
\end{tabular}
\end{center}

\begin{checkpointbox}
Du musst die Vokalzeichen nicht auswendig zeichnen können -- du musst sie
nur \emph{erkennen}. Im echten Alltag (Straßenschilder, WhatsApp)
verschwinden sie fast komplett -- Muttersprachler lesen Arabisch meistens
ohne sie, so wie du ,,Str.`` automatisch als ,,Straße`` liest.
\end{checkpointbox}
